\documentclass[10pt]{article}
\usepackage[margin=20mm]{geometry}
\usepackage{amsmath,amssymb,mathtools,fontspec,unicode-math,xurl,hyperref}
\IfFontExistsTF{DejaVu Serif}{\setmainfont{DejaVu Serif}}{\setmainfont{Latin Modern Roman}}
\setmathfont{latinmodern-math.otf}
\setlength{\parindent}{0pt}\setlength{\parskip}{6pt}\sloppy
\title{Initial Scientific Review of
Sharp Commutator Estimates of All Order
for Coulomb and Riesz Modulated Energies}
\author{Bounded scientific review: M25-0424}
\date{}
\begin{document}
\maketitle
Abstract: Eight supported problems are reported for the acquired 72-page manuscript: a counterexample to the literal finite-height localization, an incorrect transport hierarchy, two false auxiliary energy bounds, a nonmonotone truncation argument, a derivative-budget gap, false uniform amplitude dependence in a corollary, and two local algebraic or mollifier defects. The first six problems affect a stated core estimate or its derivation and are classified Major; the last two are Minor. Additional norm-power slips within Problem 6 are local Minor subfindings. These classifications do not assert failure of every principal theorem. The global first-order estimate and intended mean-field limit theorem are not refuted here. Scientific adjudication remains pending.\par
\section*{Identity, version, and scope}
The reviewed paper is by Matthew Rosenzweig and Sylvia Serfaty [source], Sharp commutator estimates of all order for Coulomb and Riesz modulated energies, DOI https://doi.org/10.1002/cpa.70010. The acquired source is arXiv:2407.15650v2, dated 28 October 2025, 72 pages. Page numbers below are the printed and physical PDF page numbers, which agree. The preserved original review records a complete reading of every page, including Appendix A and the references. This continuation reuses that documented scope and checks the eight existing problems against the bound source and retained evidence. The examined PDF has SHA256 `0b6f9013592e6d07054c32766cfb89860cbb97553c3a7a4775d2f046dfc9e70b`.\par
The arXiv record labels this version final and lists v1 and v2. The author’s publication page identifies the eventual journal publication as volume 79, issue 2 (2026), pp. 207–292 [authorrecord]. The publisher’s indexed record gives first publication on 24 September 2025 [publisherrecord]. The annual task’s 2025 identity is retained; the acquired manuscript is not silently replaced by a different journal PDF. A full publisher-page retrieval failed. Indexed publisher text corroborates the fixed-\(v\) definitions and transport identity discussed in Problem 2, but it does not establish whole-document equivalence.\par
The source itself, p. 29, footnote 4, states that an editing discrepancy in the journal’s (3.46) and subsequent computations is corrected in this version. That correction was read and visually checked. The already corrected journal discrepancy is not counted as a new problem here. Targeted checks of the arXiv, publisher, and author records did not identify a further correction addressing the findings below; this limited search does not prove that none exists.\par
\section*{Framework and checked hypotheses}
Write \[\nu_N=N^{-1}\sum_{i=1}^N\delta_{x_i}-\mu,\qquad
 K_n(x,y)=\nabla^{\otimes n}g(x-y):(v(x)-v(y))^{\otimes n}.\] The energy integrals in the statements remove the particle diagonal. The logarithmic examples use \(g=-\log|\cdot|\). The three-dimensional Coulomb example uses \(g=|\cdot|^{-1}\). Thus their normalizations agree with the source. The extension has \(k=1,\gamma=0\) for \(d=1,s=0\), and \(k=0\) for the Coulomb cases \(d=2,s=0\) and \(d=3,s=1\).\par
Theorem 1.2 assumes \(\mu\in L^1\cap L^\infty\), mass one, and absolute logarithmic interaction integrability when \(s=0\). It does not assume \(\mu\geq0\); the removal of nonnegativity is explicitly discussed on p. 17. It sets \(\lambda=(N\left\lVert \mu\right\rVert_{L^\infty(\Omega)})^{-1/d}<1\). For (1.12), \(\Omega\) contains the \(2\lambda\)-neighborhood of \(\operatorname{supp}\nabla v\), with \(v\) Lipschitz. For (1.13), the additional assumptions include \(v\in C^{2n-1,1}\), a ball \(\Omega'\) of radius \(\ell\) containing a neighborhood of \(\operatorname{supp}v\), and a \(5\ell\)-neighborhood of \(\Omega'\) in \(\Omega\). Section 5 uses a \(2\lambda\)-neighborhood instead of the statement’s \(\lambda\)-neighborhood; the examples below satisfy the stronger condition whenever Section 5 is used. All particle configurations used here are pairwise distinct, and all backgrounds in the discrete examples are bounded and compactly supported. Their absolute logarithmic integrals are therefore finite.\par
\section*{Substantive problems}
Problem 1 (Finite-height localization without compact support: Major).  The literal estimate (1.12), p. 5, is false for the stated signed-background class and a Lipschitz field whose derivative has compact support but whose values do not vanish at infinity.\par
\subsection*{Decisive example and complete verification}
Take \(d=1,s=0,N=2\), \(\Omega=[-3,3]\), and \(\ell=3\). Fix a smooth bounded field \(v\) which is \(-1\) on \((-\infty,-1]\), zero on \([-1/2,1/2]\), and \(1\) on \([1,\infty)\); its two transitions are confined to the remaining intervals. Set, for \(R>10\) and an integer \(j\geq1\), \[\begin{aligned}
 x_1&=R-\tfrac14,&x_2&=R+\tfrac14,\\
 p_R(y)&=2\,\mathbf1_{[-R-1/4,-R+1/4]}(y),&
 q_j(y)&=8\cos(2\pi j y)\,\mathbf1_{[-1/2,1/2]}(y),&
 \mu_{R,j}&=p_R+q_j .\end{aligned}\] The background has mass one and belongs to \(L^1\cap L^\infty\). Its local and global essential supremum norms are exactly 8, independently of \(R,j\). Thus \(\lambda=1/16\), the required \(2\lambda\)-neighborhood of \(\operatorname{supp}v'\) is contained in \(\Omega\), and \(\ell>\lambda\). Neither particle belongs to \(\Omega\), so \(\#I_\Omega=0\). The source radii (2.17) are both \(r_i=1/64\). On \(Q=[-3,3]\times[-3,3]\), truncation leaves their potentials unchanged.\par
Here \(q_j\) is a device for keeping the local density norm fixed; it cannot simply be discarded without taking its energy into account. Its mean is zero, its \(L^2\) norm squared is 32, and \[\widehat q_j(\xi)=8\left[
 \frac{\sin((\xi-2\pi j)/2)}{\xi-2\pi j}
 +\frac{\sin((\xi+2\pi j)/2)}{\xi+2\pi j}\right].\] For each fixed \(\xi\), this tends to zero. Uniformly in \(j\), \(|\widehat q_j(\xi)|\leq C|\xi|\) near zero, by its zero mass and bounded first absolute moment. Split \(\int |\widehat q_j(\xi)|^2|\xi|^{-1}\,\mathrm d\xi\) at a fixed \(L\). Dominated convergence controls the bounded interval, and Plancherel bounds the remaining integral by \(C/L\). Consequently this homogeneous \(H^{-1/2}\) norm tends to zero. The Fourier representation of the planar logarithmic single-layer potential then gives \[\left\lVert \nabla(g*q_j)\right\rVert_{L^2(\mathbb R^2)}\longrightarrow0 .\] This also follows directly by Fourier transformation in the horizontal coordinate. It proves convergence of the local energy on \(Q\).\par
Equation reference: eq:localcounter, eq:localenergy. Since \(v=0\) on \(\operatorname{supp}q_j\), its self-commutator vanishes exactly. All its interactions with the other two supports tend to zero by oscillatory weak convergence against smooth separated kernels. Taking \(j\to\infty\) at a fixed \(R\), the commutator in (1.12) therefore tends to \[
 I_R=2\sum_{a=R\pm1/4}\int\frac{p_R(y)}{a-y}\,\mathrm dy
     =4\log\frac{2R+1/2}{2R-1/2}\sim\frac2R .\] Let \(h_R=g*(\tfrac12\delta_{x_1}+\tfrac12\delta_{x_2}-p_R)\). Every source point is horizontally at least \(R-13/4\) from \(Q\). The total variation of this source is 2 and \(|\nabla g(X)|=|X|^{-1}\); hence \[
 \int_Q|\nabla h_R|^2\leq\frac{144}{(R-13/4)^2}.\] If (1.12) held with its stated universal constant, the \(j\)-limit would give \(I_R\leq C\left\lVert v'\right\rVert_\infty\,144/(R-13/4)^2\) for every \(R\). The left side is of order \(R^{-1}\) and the right side is of order \(R^{-2}\), an impossibility. This is an analytic counterexample sequence, not an inference from finite numerical samples.\par
\subsection*{Origin, alternate readings, and impact}
Remark 3.1, p. 23, uses \(\widetilde v(x,z)=\chi(z/\ell)(v(x),0)\). Its vertical derivative contains \(\ell^{-1}\chi'(z/\ell)v(x)\), which remains nonzero on a nonzero plateau where \(v'(x)=0\). Thus the support replacement invoked in the local arguments on pp. 27 and 31 is unavailable for this field. The compact support premise needed for (3.5) is absent from the first part of Theorem 1.2.\par
A natural repair route is to retain the full vertical domain \(\Omega\times\mathbb R^k\) with the trivial extension. Another is to impose the compact-support and domain hypotheses needed by the cutoff proof and use the actual support of \(\nabla\widetilde v\). A full replacement theorem under either amended set of assumptions is not certified here. A nonnegative-background-only statement would also exclude this particular counterexample, but that is a genuine restriction of the printed class. A fixed finite change to \(\ell\) does not repair the \(R^{-1}\) versus \(R^{-2}\) discrepancy. The example does not refute the compactly supported higher-order estimate, the global estimate (1.7), or the mean-field application of that global estimate.\par
Problem 2 (The transported commutator hierarchy omits a chain-rule term: Major).  Equations (1.31)–(1.32), p. 8, and Lemma A.3, especially (A.6)–(A.10), p. 63, use an incorrect transport identity for the explicitly fixed field \(v\).\par
\subsection*{Exact differentiation}
Equation reference: eq:correcttransport. For the source definition \[\kappa_t^m(x)=\int\nabla^{\otimes m}g(x-y-tv(y)):
                        (v(x)-v(y))^{\otimes m}\,\mathrm df(y),\] differentiation gives, for \(m\geq1\), \[\begin{aligned}

 (\partial_t+v\cdot\nabla)\kappa_t^m
  ={}&\kappa_t^{m+1}\\
 &+m\int\nabla^{\otimes m}g(x-y-tv(y)):
 \bigl((v\cdot\nabla v)(x)\otimes
       (v(x)-v(y))^{\otimes(m-1)}\bigr)\,\mathrm df(y).\nonumber\end{aligned}\] The last term differentiates the factors \(v(x)-v(y)\). It is absent from (A.7). For \(m=1\), it is exactly \((v\cdot\nabla v)(x)\cdot\nabla h_t(x)\).\par
This defect occurs with all the qualitative conventions used in the source. Take \(d=3,s=1\), and a smooth compactly supported \(v\) equal to \(x\) on \(B(0,3)\). Let \(\widehat\psi\) be an even, nonnegative smooth function supported in \(1<|\xi|<2\) with \(\int \xi_1^2\widehat\psi(\xi)\,\mathrm d\xi>0\). Choose \(a>0\) so that \(h(x)=a\,\partial_1\psi(x-e_1)\) satisfies \(\partial_1h(e_1)=-2\), and set \(f=-\Delta h/(4\pi)\). Then \(f\) is Schwartz, has zero mean, has Fourier support away from zero, and \(g*f=h\), since \(-\Delta g=4\pi\delta_0\). At \(x=e_1,t=0\), the omitted term is exactly \(-2\). Thus even the standing Fourier-gap convention after (3.21), p. 25, does not eliminate the counterexample.\par
For transparent independent arithmetic, one may also take \(f=\phi-\phi(\,\cdot-2e_1)\), with \(\phi\) a radial smooth unit-mass bump in \(B(0,1/4)\). Newton’s identity gives, near \(x=e_1\) on the axis, \[\kappa_t^1(xe_1)=-\frac1x+\frac{2-x}{(2+2t-x)^2}.\] At \(x=1,t=0\), direct differentiation gives \(\partial_t\kappa_t^1=-4\), \(v\cdot\nabla\kappa_t^1=2\), and \(\kappa_0^2=0\). The discrepancy is again \(-2\). This second example illustrates the calculation; the preceding annular Fourier construction meets the additional Fourier-support convention exactly.\par
\subsection*{Dependency and repair boundary}
Lemma A.5 uses (A.7) in (A.16), p. 64; Lemma A.7 and the proof of Proposition A.1 on pp. 67–69 then use that representation. Therefore the advertised alternative proof is not valid as written. Introducing the transported velocity \(v_t=v\circ(I+tv)^{-1}\), which satisfies a Burgers transport relation, can cancel a corresponding acceleration term in a different hierarchy. It does not validate (A.6) with its fixed \(v\); definitions and the coefficient induction would have to be redone. Section 4 supplies a separate PDE approach, so this finding does not by itself falsify the estimate in Proposition A.1 or the principal estimates.\par
Problem 3 (Two false auxiliary energy bounds in the local renormalization: Major).  The integral in Lemma 5.1, (5.7), p. 51, can contain an unsmeared particle singularity. Independently, the global neutral-cell bound (5.40), p. 55, fails for the logarithmic kernel even for nonnegative backgrounds.\par
\subsection*{An infinite left side in Lemma 5.1}
Section 5 sets \(\eta_i=r_i/2\) for \(x_i\in\Omega'\) and \(\eta_i=0\) otherwise, explicitly leaving the other particles as Dirac masses ((5.1), p. 50). Yet (5.7) integrates \(|z|^\gamma|\nabla\bar h_{N,\vec\eta}|^2\) over \((\operatorname{supp}\nabla\widetilde v)\cup(\Omega\times[-\lambda,\lambda]^k)\).\par
Take \(d=1,s=0,N=2\), \(\mu=\mathbf1_{[-1/2,1/2]}\), particles \(0,4\), a smooth \(v\) supported in \([-1,1]\), \(\Omega'=[-3,3]\), and \(\Omega=[-18,18]\). Here \(\ell=3,\lambda=1/2\); all the stronger Section 5 neighborhood conditions hold. The particle at 4 lies in \(\Omega\setminus\Omega'\) and is not smeared. Near \(X=(4,0)\), \[\bar h_{N,\vec\eta}(X)=-\tfrac12\log|X-(4,0)|+H(X),
 \qquad H\ \hbox{smooth there}.\] The energy on an annulus of radii \(\varepsilon,\varepsilon_0\) is \((\pi/2)\log(\varepsilon_0/\varepsilon)+O(1)\), and diverges. The right side \(\Xi\) in (5.2) is finite: it uses \(h_{N,\vec r}\), which truncates both particles at positive radii. The closing appeal to (2.30) on p. 52 cannot remedy the unsmeared point.\par
\subsection*{A separated logarithmic cell with unbounded global energy}
Now take \(d=2,s=0,N=2\), particles \(0,e_1\), \(\Omega'=\overline{B(0,3)}\), \(\Omega=\overline{B(0,18)}\), and a fixed smooth \(v\) supported in \(B(0,1)\). Let \(\nu_b\) be the uniform probability density on the disk of radius \(1/10\) centered at \(b\), and, for \(R>50\), set \[\mu_R=\tfrac14\nu_0+\tfrac14\nu_{Re_1}
                         +\tfrac12\nu_{3Re_1}.\] This is a nonnegative bounded probability density. In \(\Omega\) its essential supremum is \(25/\pi\), so \(\lambda=\sqrt{\pi/50}<1\), independently of \(R\). The radii \(\eta_i=\lambda/8\) are fixed. Choose a fixed radial smooth unit-mass mollifier at this scale with the required vanishing moments. Such a signed mollifier can be obtained from two radial dilates with coefficients \(4/3,-1/3\), which cancel the second moments. Its support can be made smaller than \(\eta_i/2\).\par
Equation reference: eq:logcell. For the particle at zero, \(R_i=2R\) satisfies the defining cell condition \(\int_{B(0,R_i)}\mu_R=1/N\) used in (5.24). Its neutral source is \[f_R=\tfrac12\rho^{(\eta_i)}
       -\mu_R\mathbf1_{B(0,2R)}
     =a_0-\tfrac14\nu_{Re_1},\qquad
 a_0=\tfrac12\rho^{(\eta_i)}-\tfrac14\nu_0,\quad \int a_0=\tfrac14 .\] By the radial mean-value property of the logarithmic kernel on disjoint disks, the cross interaction is exactly a mass product times \(-\log R\). Consequently, with \(E_0\) independent of \(R\), \[
 \iint-\log|x-y|\,\mathrm df_R(x)\,\mathrm df_R(y)=E_0+\tfrac18\log R,\qquad
 \int_{\mathbb R^2}|\nabla(g*f_R)|^2=2\pi E_0+\tfrac\pi4\log R .\] The \(p=0\) instance of (5.40) instead bounds this global energy by \(C/(N^2\eta_i^s)=C/4\), a contradiction. The intermediate assertion on p. 55 that the logarithmic potential of a unit-mass mollifier is uniformly bounded also fails: it equals \(-\log|x|+O(1)\) at infinity.\par
\subsection*{Other hypotheses, dependency, and bounded repair}
The two examples are independent and use nonnegative \(\mu\). The printed theorem’s larger signed class presents an additional issue: \(\int_{B(x_i,R_i)}\mu=1/N\) does not bound \(\int_{B(x_i,R_i)}|\mu|\), as needed when estimating the integrals in (5.28)–(5.39). Signed moment-cancelling mollifiers likewise require total-variation constants instead of a probability-density argument.\par
Lemma 5.1 is used in the estimate of the smeared term and in the passage after (5.45), p. 56. The neutral-cell estimate enters the local regularity step for Term 2. One must restrict the energy to regions without unsmeared singularities and retain a genuinely local cell norm, or supply another valid estimate. Merely enlarging a local norm to the global energy loses the logarithmic cancellation illustrated above. These counterexamples invalidate the displayed auxiliary bounds and their current uses; they do not prove that the intended high-order commutator estimate is false. Section 5.4’s global treatment uses a different, cruder Term 2 estimate and is not refuted by [eq:logcell] (\#eq:logcell) alone.\par
Problem 4 (The mesoscale truncation path is not monotone: Major).  The monotonicity used between (2.46) and (2.47), p. 21, does not follow from Lemma 2.2. It is false for the path (2.44), even after correcting the fixed-potential typo in (2.45).\par
\subsection*{Exact counterevidence to the proof step}
The path chooses \(\eta_i^k=2^k\lambda\) while \(\operatorname{dist}(x_i,\partial\Omega)\geq2^k\lambda\), and otherwise resets it to \(r_i\). A point at distance \(3\lambda\), with \(r_i=\lambda/4\), therefore has the successive radii \[\lambda,\quad2\lambda,\quad\lambda/4 .\] This is not a coordinatewise increasing path.\par
The issue is not only logical. In \(d=3,s=1\), take the half-space \(\Omega=\{x_3\geq0\}\), particles \(x_1=(0,0,3\lambda)\), \(x_2=(3\lambda,0,3\lambda)\), and \(0<\lambda<1\). Put a bounded probability density of height \(1/(2\lambda^3)\) in a cube of volume \(2\lambda^3\), far from these particles and their balls but inside \(\Omega\). This makes the prescribed local scale exactly \(\lambda\), and \(r_i=\lambda/4\). All balls of radius \(2\lambda\) stay in \(\Omega\).\par
Equation reference: eq:shell. For two uniform Coulomb shells of common radius \(a=2\lambda\), separated by \(r=3\lambda\), their mutual energy is \[
 J=\frac{1}{2ar}\int_{|r-a|}^{r+a}
           u\min\{u^{-1},a^{-1}\}\,\mathrm du
   =\frac{5}{16\lambda}<\frac{1}{3\lambda}.\] At radii \(\lambda\) or \(\lambda/4\), the shells are disjoint, and their mutual energy equals the point interaction \(1/(3\lambda)\). In the renormalized \(\mathcal F^{\vec\eta}\), self terms cancel by its definition and the background contributes no change. Thus, for \(N=2\), \[\mathcal F^{\vec\eta^0}-\mathcal F^{\vec\eta^1}
 =\frac1{192\lambda},\qquad
 \mathcal F^{\vec\eta^2}-\mathcal F^{\vec\eta^1}
 =\frac1{192\lambda}>0 .\] The asserted nonincrease in \(k\) is disproved. This example targets the proof’s universal monotonicity assertion; the selected-pair sum in (2.40) need not be nonempty, so it is not a counterexample to (2.40).\par
As printed, (2.45) also retains \(h_{N,\vec r}\) instead of \(h_{N,\vec\eta}\). Its energy term then does not vary with \(\vec\eta\). With the background outside all truncation balls, even increasing the radii makes the negative self-energy term increase. Replacing this potential is necessary but leaves the reset obstruction above. For \(k=1\), the preceding vertical-domain replacement additionally uses coincidence only above height \(2^{k+2}\lambda\), which can exceed \(\ell\); the asserted replacement by a slab of height \(\ell\) requires a separate justification.\par
\subsection*{Dependency and repair boundary}
Proposition 2.5, pp. 20–23, is invoked in (5.14), (5.57), and the later macroscopic estimates. A proof must use admissible comparisons on a fixed interior set, or separately estimate the reset terms and the vertical tails. No completed replacement argument is supplied here. There is also a local sign error in (2.53): for \(f/g=t^{-a}\), its derivative is \(-at^{-a-1}\), which is at most, not at least, \(-a(2^k\lambda)^{-a-1}\) on the interval. For \(\lambda=1/10,a=1,k=2,t=1/10\), the two values are \(-100\) and \(-25/4\). That particular inequality can be repaired with a factor \(4^{a+1}\); it does not repair the nonmonotone path.\par
Problem 5 (The stated high-order derivative budget is not obtained: Major).  Theorem 1.2, (1.13), p. 5, has undefined summation indices and a regularity budget not delivered by the estimates concluding its proof.\par
\subsection*{Index and derivative calculation}
The power of \(\lambda\) printed in (1.13) is \(\lambda^{-(n-p)+\sum_{k=1}^p c_{n-k}}\), while the only bound indices are \(c_1,\ldots,c_{n-p}\). At \(n=3,p=2\), the exponent uses \(c_2\), although only \(c_1\) is bound. At \(p=0\), the exponent loses every derivative order. This is not a defined general formula.\par
Correcting only that index does not settle the proof. Theorem 4.9, (4.119), p. 47, feeds into (5.34), p. 54, and (5.42), p. 55: \[A_{v,p}=\sum_{\substack{a_1,\ldots,a_{n-p}\geq0\\
                        \sum a_i\leq2n}}
 \lambda^{\sum a_i}\prod_{i=1}^{n-p}\left\lVert D^{1+a_i}v\right\rVert_\infty .\] The exact change \(b_i=1+a_i\) gives \[n-p\leq\sum b_i\leq3n-p,\qquad
 \lambda^{\sum a_i}=\lambda^{\sum b_i-(n-p)},\] whereas (1.13) has the upper bound \(2n\). In particular \(p=n-1,a_1=2n\) uses \(D^{2n+1}v\). The hypothesis \(C^{2n-1,1}\) controls derivatives through order \(2n\), not \(2n+1\).\par
This is a real loss of control, not an assertion that high derivatives must happen to be large for every field. For a fixed smooth compact cutoff \(\chi\) with a plateau, take \[v_J=v_0+J^{-2n}\chi(x)\sin(Jx)\] in one dimension, where \(v_0''\) is nonzero on that plateau. The \(C^{2n-1,1}\) norms stay bounded and the supports stay fixed, while \(\left\lVert D^{2n+1}v_J\right\rVert_\infty\) grows like \(J\). The multiplier \((\ell\left\lVert D^2v_J\right\rVert_\infty)^{n-1}\) of the relevant term in (5.47)/(5.61) does not tend to zero. For fixed admissible \(\lambda,\ell\), the resulting bound cannot be controlled by the printed derivative data.\par
\subsection*{Impact and possible repair}
The concluding assertion after (5.61), p. 58, therefore needs a sharper regularity argument or a corrected weaker theorem with the derivative budget actually proved. The convention \(A_{v,0}:=1\) printed after (5.42) also has the wrong endpoint: the empty product is at \(p=n\). These observations do not constitute a counterexample to an appropriately defined optimal high-order estimate. The Major classification concerns the missing derivation of the announced regularity and scale dependence.\par
Problem 6 (Uniform amplitude dependence in Corollary 1.3: Major).  Corollary 1.3, (1.14), p. 6, has \(CM\ell^{-n}\) with the inherited constant independent of \(M\). Its left side is homogeneous of degree \(n\) in the field, so \(M^n\) is required. Related squared-norm displays also lose a power.\par
\subsection*{A nonzero example, rather than homogeneity alone}
Take \(d=1,s=0,N=2\), \(\mu=\mathbf1_{[-1/2,1/2]}\), and particles \(-1/4,1/4\). Choose \(v_0\) smooth, equal to \(x\) on \([-2,2]\), and supported in \([-4,4]\), with \(\Omega'=[-6,6]\), \(\Omega=[-36,36]\), and \(\ell=6\). Here \(\lambda=1/2\), and even the stronger Section 5 neighborhood conditions hold. Put \(v_A=A v_0\) for \(A>0\). On all the relevant background and particle supports, \[K_n(x,y)=(-1)^n(n-1)!\,A^n,\qquad x\ne y.\] The total mass of \(\nu_N^{\otimes2}\) off the diagonal is \(-1/N\). In particular the signed \(n=2\) integral equals \(-A^2/2\), and its absolute value is \(A^2/2\). The energy bracket in (1.14) is finite and independent of \(A\). One can choose \(M=A M_0\), where \(M_0=\max_{0\leq k\leq4}\ell^k\left\lVert D^kv_0\right\rVert_\infty\). A bound linear in \(M\) cannot hold as \(A\to\infty\).\par
The source explicitly presents Corollary 1.3 as the precise scale and homogeneity consequence of (1.13), p. 5. The example therefore contradicts a stated quantitative guarantee tied to the central local estimate. This corollary component is Major under the core-contribution criterion. Replacing \(M\) by \(M^n\) restores the required homogeneity, but does not prove the other estimates needed for the corollary. Allowing the constant to depend on \(M\) would change its stated uniformity. The simplicity of the normalization repair does not remove this dependency.\par
\subsection*{Local norm-power subfindings: Minor}
The same scaling check shows that (1.26), p. 7, needs \(\left\lVert \nabla v\right\rVert_\infty^2\), because it bounds a squared gradient norm; (3.17), p. 25, has the consistent unsquared norm and is counterevidence against treating this as a failure of the whole first-order argument. Similarly the coefficient in the squared high-order norm (1.33), p. 8, must have degree \(2n\) in \(v\), rather than \(n\). Finally, the duality test stated on p. 27 for (3.19) requires \((-\Delta)^{(d-s)/2}\kappa\): using \((-\Delta)^{d-s}\kappa\) gives \(\int\kappa(-\Delta)^{d-s}\kappa=\left\lVert \kappa\right\rVert_{\dot H^{d-s}}^2\), not the claimed \(\dot H^{(d-s)/2}\) square. These are local power corrections with explicit consistent alternatives.\par
Problem 7 (The Term 3 decomposition omits the diagonal self correction: Minor).  The reduction to Term 2 plus the off-diagonal Term 4 in (5.48), p. 56, omits the smooth–atomic self interactions which survive removal of the particle diagonal.\par
\subsection*{Exact algebra and example}
Equation reference: eq:diagonal. Abbreviate \(\rho_i=\rho_{x_i}^{(\eta_i)}\) for the smeared particles. Using the definitions (5.17)–(5.19), p. 52, direct expansion gives \[\begin{aligned}

 \mathrm{Term}_3-\mathrm{Term}_2
 ={}&\mathrm{Term}_4\\
 &+\frac1{N^2}\sum_{i\in I_{\Omega'}}
 \left(\iint K_n(x,y)\rho_i(x)\rho_i(y)\,\mathrm dx\,\mathrm dy
       -2\int K_n(x_i,y)\rho_i(y)\,\mathrm dy\right).\nonumber\end{aligned}\] Only the Dirac–Dirac diagonal is removed; the other two kinds of self interaction are not removed by that operation.\par
Use the \(n=2,A=1\) example in Problem 6 and unit-mass mollifiers with the source’s small supports. Then \(K_2=1\) on every interacting support. Both the fully smoothed signed measure and each \(\delta_{x_i}-\rho_i\) have mass zero, so \[\mathrm{Term}_1=\mathrm{Term}_2=\mathrm{Term}_4=0,\qquad
 \mathrm{Term}_3=-1/N=-1/2 .\] The missing sum in [eq:diagonal] (\#eq:diagonal) is exactly \(-1/2\), even when the moment-cancelling mollifiers are signed.\par
\subsection*{Repair and limited impact}
For fixed smooth mollifier shapes, \(|K_n(x,y)|\leq
C\left\lVert \nabla v\right\rVert_\infty^n|x-y|^{-s}\), with the corresponding bounded kernel estimate when \(s=0\). Scaling the integrals in [eq:diagonal] (\#eq:diagonal), using \(s<d\), bounds each by \(C\left\lVert \nabla v\right\rVert_\infty^n/(N^2\eta_i^s)\). For \(s>0\), the self-energy budget (2.28) is the appropriate place to absorb their sum. For \(s=0\), the density-error term in \(\Xi\) directly controls \(\#I_{\Omega'}/N^2\), since \(\lambda^d\left\lVert \mu\right\rVert_{L^\infty(\widehat\Omega)}\geq1/N\). Total-variation constants of signed mollifiers must be retained. Thus the equality is false as written, but a local error term of the expected size is available. This finding alone does not falsify the claimed final inequality.\par
Problem 8 (The proposed moment-cancelling mollifier has the wrong jet: Minor).  Formula (5.3), p. 50, does not give the stated vanishing moments for arbitrary smooth unit-mass bumps \(\psi,\zeta\).\par
\subsection*{Second-order check and constructive repair}
Use the Fourier convention \(\widehat f(\xi)=\int e^{-i\xi x}f(x)\,\mathrm dx\), in one dimension with \(m=2\). Let \(A,B\) be the first moments of \(\psi,\zeta\), and \(S,T\) their second moments. Their Taylor series are \[\widehat\psi=1-iA\xi-\tfrac12S\xi^2+O(\xi^3),\qquad
 \widehat\zeta=1-iB\xi-\tfrac12T\xi^2+O(\xi^3).\] The formula subtracts \(\widehat\zeta(-iA\xi-S\xi^2/2)\). It therefore gives \[\widehat\rho(\xi)=1+AB\xi^2+O(\xi^3),\qquad
 \int x^2\rho(x)\,\mathrm dx=-2AB .\] Take both bumps to be translates of an even smooth probability bump by the displacement 1; then \(A=B=1\), and the purportedly zero second moment is \(-2\).\par
For every order \(m\), a valid repair is to set \[P(\xi)=\operatorname{Taylor}_m
           \left(\frac{\widehat\psi(\xi)-1}{\widehat\zeta(\xi)}\right),
 \qquad \widehat\rho=\widehat\psi-\widehat\zeta P .\] The denominator is nonzero near zero, \(P(0)=0\), and \(\widehat\rho=1+O(|\xi|^{m+1})\). In several dimensions use the total-degree Taylor polynomial. Fourier inversion gives a finite combination of derivatives of compactly supported smooth bumps, hence the required compact support and smoothness. For the actually used \(m\leq2\), centered symmetric choices also avoid the specific second-order cross term.\par
The moments are used in the Taylor estimate (5.51), p. 57, and \(m=2\) is selected in the Coulomb case on p. 58. Their construction is therefore operative, but this repair is local. A smooth unit-mass function with all second moments zero cannot be nonnegative: otherwise its mass would be supported at the origin. Estimates using the repaired mollifier must accordingly allow a signed density and keep its fixed total-variation constants. The finding does not show a loss of the intended rate once a valid mollifier and valid surrounding estimates are supplied.\par
\section*{Dependency assessment and conclusion}
The distinction between a false displayed assertion and an unproved intended conclusion is essential here. Problem 1 disproves (1.12) in its literal signed-background, noncompact-field class. Problem 6 disproves the printed uniform \(M\)-dependence of the corollary but has an elementary normalization repair. Problems 2–5 demonstrate failures in substantial derivations or auxiliary assertions; they do not supply counterexamples to all possible correct formulations of the principal high-order estimates. Problems 7–8 have explicit local algebraic repairs.\par
The primary reading followed the first-order stress identity and truncation arguments, the Section 4 commutator PDE and regularity inductions, the Section 5 renormalization and global/macroscopic variants, the mean-field application, and Appendix A’s alternative induction. Singularity, signedness, diagonal removal, finite-height localization, low-frequency conventions, and derivative counts were checked where operative. No independent counterexample to the global first-order result or the intended mean-field theorem was established. Nor is this report a certification that every uncriticized argument is correct. The cited prior literature was not all independently re-proved or fully retrieved; the decisive evidence here is self-contained.\par
Recorded dimension | State\par
Supported problems | 6 Major; 2 Minor\par
Scientific decision | Pending\par
Human reviews | None recorded\par
Independent acceptance | Not conferred by this initial report\par
Author adoption of suggested repairs | Not established\par
\section*{Independent checks and evidence boundaries}
The original reviewer wrote and executed `verification/0424/independent\_checks/independent\_checks\_001.py` using the existing Python runtime, SymPy 1.13.3, and mpmath 1.3.0. It checks the transport derivatives, shell integral, derivative-index arithmetic, logarithmic affine kernels, and mollifier jets, and records numerical illustrations of [eq:localcounter] (\#eq:localcounter)–[eq:localenergy] (\#eq:localenergy). The process exited with code zero and all its assertions passed. The results file has SHA256 `7d853ae94d53dc34edd65a80ec9ee2c57d524457f9b819497b625330bc80c057`. Floating-point illustrations are not interval certificates or substitutes for the analytic limiting arguments. No author code was executed.\par
The source PDF and manifests, exact TeX, genuinely compiled PDF, compiler and dependency payloads, logs, exact-PDF parsing, render hashes, and actual page-view declarations are retained under this paper’s own verification tree. The final support inventory is published after these artifacts. Rendering or transmitting an image is not counted as viewing it. The source text was fully read; the evidence pages cited in the problems were also inspected through the authorized native image route. All pages of the delivered report must pass that same actual inspection before it is marked ready.\par
The preserved original review supplies the full reading, calculations, and source-image inspection records. Recovery continuation rechecks the existing claims with native bounded proof/source helpers, corrects the classification of Problem 6, and completes the current build, page views, and metadata. Historical executions and newly performed work are attributed separately in provenance; a current check is not attributed retrospectively. The supplied UI selection and the runtime’s self-description are also recorded separately, without a backend-equivalence claim. This is an initial AI-assisted review, not a recorded human review or an independent acceptance decision.\par
9 M. Rosenzweig and S. Serfaty, Sharp commutator estimates of all order for Coulomb and Riesz modulated energies, arXiv:2407.15650v2 (28 October 2025), 72 pp. https://arxiv.org/abs/2407.15650. DOI: https://doi.org/10.1002/cpa.70010. The exact acquired PDF is the scientific source for all numbered locators. S. Serfaty, author publication record, https://math.nyu.edu/\textasciitilde{}serfaty/papers.html, consulted 29 October 2026 for publication and correction context. The retrieved page is retained in `verification/0424/web\_checks/0424\_web\_primary.json`. This record is not used as a substitute for mathematical evidence. Wiley, indexed publisher record for DOI 10.1002/cpa.70010, https://onlinelibrary.wiley.com/doi/10.1002/cpa.70010, consulted 29 October 2026. The full-page retrieval failed; only the preserved indexed passages and publication metadata were used.\par
\end{document}
